\section{Conclusion}
\label{sec:conclusion}

We have developed a generalized hierarchical Bayesian method for multiple-changepoint segmentation with irregular designs, multiple aligned sequences, and known or latent group structure. For regular exponential-family segment models with proper conjugate priors, sufficient statistics give exact segment marginal likelihoods and posterior moments. These quantities enter common dynamic-programming recursions for the posterior distribution over segment counts and changepoint locations. A separate max-sum recursion with backtracking returns the joint MAP partition, so marginal posterior summaries and joint optimization are not conflated.

The same partition formulation accommodates design-dependent segment cohesion and boundary hazards, exact pooling of conditionally independent sequences with common changepoints, and a finite latent-group criterion optimized by alternating minorization and maximization. Nonconjugate segment models can be incorporated through numerical marginal-likelihood calculations; the resulting posterior approximation is related to the intended posterior only when a segmentwise error bound such as Assumption~\ref{ass:block-error} is established. Posterior prediction is likewise defined by the fitted observation family and its coordinate support.

The numerical studies in Section~\ref{sec:results} cover posterior calibration across 22 preregistered cells, observation-family generality across five exponential families, design-aware partition priors, exact shared-boundary pooling, finite-template clustering, nonconjugate block-evidence approximations, robustness to misspecification and missing data, posterior prediction, computational scaling, and three applied datasets (well-log, array-CGH, and methylation). In the reported clustering configuration, hard allocation accuracy is 96\%; this characterizes the specified simulation. Timing results describe the tested implementation over the measured sequence lengths. For the applied datasets, boundaries are reported as posterior summaries of the fitted model; no independently annotated reference was available for accuracy evaluation against ground truth.

The statistical contribution is the integration of conjugate segment models, exact posterior calculations for locally factorized partitions, irregular-coordinate priors, shared-boundary pooling, and a finite-template clustering criterion under explicitly stated assumptions. Together these provide a principled, computationally tractable framework for hierarchical Bayesian segmentation across diverse observation families and experimental designs.
